\section{Evaluation}
\label{sec:evaluation}

In this section, we evaluate our approach using 62 real-world C programs.
%
We first describe our implementation of Forcrat, which realizes the proposed
approach (\cref{sec:evaluation:implementation}), the benchmark programs used
for evaluation (\cref{sec:evaluation:benchmark}), and the overall experimental
process (\cref{sec:evaluation:process}).
%
The implementation and benchmarks are publicly available~\cite{artifact}.
%
We then address the following research questions:
%
\begin{itemize}
  %
  \item \textbf{RQ1. Correctness}: Does it transform programs while preserving their
    semantics? (\cref{sec:evaluation:correctness})
  %
  \item \textbf{RQ2. Efficiency}: Does it analyze and transform programs efficiently?
    (\cref{sec:evaluation:efficiency})
  %
  \item \textbf{RQ3. Applicability}: Is it applicable to replacing a wide range of libc
    I/O API calls? (\cref{sec:evaluation:applicability})
  %
  \item \textbf{RQ4. Impact on performance}: What is the effect of the transformation on
    program performance? (\cref{sec:evaluation:performance})
  %
  \item \textbf{RQ5. Security}: How much does the transformation improve the
    security of programs? (\cref{sec:evaluation:security})
  %
  \item \textbf{RQ6. Code changes}: How much does the code change due to the
    transformation? (\cref{sec:evaluation:code})
  %
\end{itemize}
%
Our experiments were conducted on an Ubuntu machine with an Intel Core i7-6700K
(4 cores, 8 threads, 4GHz) and 32GB of DRAM.
%
Finally, we discuss threats to validity (\cref{sec:evaluation:threats}).

\subsection{Implementation}
\label{sec:evaluation:implementation}

We built Forcrat on top of the Rust compiler~\cite{rustc}.
%
For static analysis, we utilize Rust's mid-level intermediate representation
(MIR)~\cite{mir}, which expresses functions as control flow graphs composed of
basic blocks.
%
For program transformation, we modify the abstract syntax tree and then
pretty-print it back to code.
%
We used C2Rust v0.18.0 with minor modifications.

\subsection{Benchmark Programs}
\label{sec:evaluation:benchmark}

\renewcommand{\arraystretch}{0.45}
\begin{table}[t]
  \caption{Benchmark programs}
  \label{tab:benchmark}
  \centering
\scriptsize
  \vspace{-0.5em}
  \begin{tabular}{@{~}l@{~}||@{~}r@{~}|@{~}r@{~}|@{~}r@{~}|@{~}r@{~}}
    \textbf{Name} & \textbf{C LOC} & \textbf{Rust LOC} & \textbf{API Calls} & \textbf{Remaining Calls} \\
    \hline
avl&101&114&1&0\\
bc-1.07.1&10,810&16,982&158&36\\
brotli-1.0.9&13,173&127,691&99&0\\
bst&65&89&1&0\\
buffer-0.4.0&395&1,137&9&0\\
bzip2&5,316&13,731&226&87\\
cflow-1.7&20,601&26,375&150&11\\
compton&8,748&14,084&130&3\\
cpio-2.14&35,934&80,929&163&2\\
dap-3.10&22,420&43,549&1,391&27\\
diffutils-3.10&59,377&95,835&236&99\\
ed-1.19&2,439&5,636&66&6\\
enscript-1.6.6&34,868&78,749&987&234\\
findutils-4.9.0&80,015&139,859&311&38\\
gawk-5.2.2&58,111&140,566&643&517\\
genann-1.0.0&608&1,818&69&0\\
glpk-5.0&71,805&145,738&82&28\\
gprolog-1.5.0&52,193&74,381&583&80\\
grep-3.11&64,084&84,902&70&11\\
gsl-2.7.1&227,199&422,756&880&0\\
gzip-1.12&20,875&21,605&97&0\\
hello-2.12.1&8,340&10,688&33&11\\
heman&7,048&14,690&12&0\\
hiredis&7,305&14,042&481&0\\
indent-2.2.13&19,255&15,581&125&29\\
less-633&20,063&45,685&64&18\\
libcsv&965&3,010&9&0\\
libosip2-5.3.1&15,772&36,286&134&0\\
libtool-2.4.7&3,769&5,701&8&7\\
libtree-3.1.1&1,412&2,632&149&0\\
libzahl-1.0&2,438&4,096&1&0\\
lil&2,934&5,558&28&0\\
lodepng&5,098&14,299&11&0\\
make-4.4.1&28,911&36,336&362&1\\
mcsim-6.2.0&18,527&36,454&906&78\\
minilisp&722&2,149&24&7\\
mtools-4.0.43&18,266&37,021&627&10\\
nano-7.2&42,999&74,995&102&25\\
nettle-3.9&61,835&82,742&290&0\\
patch-2.7.6&28,215&103,839&207&41\\
pexec-1.0rc8&5,357&12,301&140&0\\
pocketlang&14,267&41,439&82&0\\
pth-2.0.7&7,590&12,950&178&9\\
quadtree-0.1.0&365&1,057&14&0\\
raygui&1,588&17,218&38&0\\
rcs-5.10.1&28,286&36,267&106&99\\
rgba&396&2,128&2&0\\
robotfindskitten&398&1,508&3&0\\
screen-4.9.0&39,335&72,201&113&3\\
sed-4.9&48,190&68,465&426&266\\
shairport&4,995&10,118&66&3\\
tar-1.34&66,172&134,970&305&80\\
time-1.9&2,828&1,830&102&0\\
tinyproxy&5,667&12,825&56&0\\
tulipindicators&12,371&22,864&45&0\\
twemproxy&22,738&74,593&199&0\\
units-2.22&7,240&11,521&241&44\\
urlparser&56&1,360&12&0\\
uucp-1.07&51,123&77,872&1,150&297\\
webdis&14,369&29,474&24&5\\
wget-1.21.4&81,188&192,742&325&125\\
which-2.21&2,010&2,241&42&0\\\hline
{\bf Total}&&&13,594&2,337\\
  \end{tabular}
  \vspace{-1em}
\end{table}

We used benchmark programs from previous studies on improving C2Rust-generated
code~\cite{zhang2023ownership, hong2024dont, hong2024tag}.
%
This resulted in 62 programs, including 36 GNU packages as well as open-source
projects from GitHub.
%
Columns 2--4 of \cref{tab:benchmark} present the C LOC, Rust LOC, and the number
of libc I/O API calls in each program, respectively.
%
The benchmarks are diverse in size, ranging up to 422,756 Rust LOC and 1,391 API
calls.
%
Among these, 32 programs are accompanied by test suites.

\subsection{Experimental Process}
\label{sec:evaluation:process}

We first translated all benchmark programs using C2Rust.
%
We checked whether the translated programs were compilable and found that 25
programs were not, due to type errors.
%
We manually fixed these errors, modifying an average of 11.7 lines per program.
%
We also ran the test suites, when available, and confirmed that all tests
passed.
%
Next, we applied our approach to the translated code, performing static analysis
and code transformation while measuring the times to answer RQ2.
%
After this, we checked whether the transformed programs were compilable and
passed their test suites to answer RQ1.
%
We also measured the runtime of each test suite to answer RQ4.
%
Finally, we inspected the transformed code to count the number of transformed
API calls and to measure the amount of unsafe code and code changes, addressing
RQ3, RQ5, and RQ6.

\subsection{RQ1: Correctness}
\label{sec:evaluation:correctness}

We evaluate the correctness of the proposed approach by checking whether the
transformed program is compilable and passes its test cases.
%
First, the experimental results show that our approach produces compilable code
for most programs.
%
Among the 62 programs, only one program, nano-7.2, fails to compile.
%
We manually investigated this case and confirmed that the failure is not due to
a problem in our approach.
%
Below is the problematic part in the original C code:

\begin{lstlisting}[language=C]
if (original != NULL) { v = copy_file(original); }
if (original == NULL || v < 0) { .. }
\end{lstlisting}

\noindent
%
Here, \C{original} is a stream, and \C{copy\_file} closes it.
%
Although the program uses \C{original} even after it is closed, it is only to
check nullity, not to perform I/O, which is not considered a bug.
%
However, Forcrat translates this into the following uncompilable Rust code:

\begin{lstlisting}[language=Rust]
if !original.is_none() { v = copy_file(original); }
if original.is_none() || v < 0 { .. }
\end{lstlisting}

\noindent
%
Since Rust's ownership type system prevents the use of a variable after it is
passed to another function, using \C{original} after calling \C{copy\_file}
results in a type error.
%
To obtain compilable code through translation, we can fix the original program
as shown below, which checks nullity and stores the result in a variable before
calling \C{copy\_file}:

\begin{lstlisting}[language=C]
int is_original_null = original == NULL;
if (!is_original_null) { v = copy_file(original); }
if (is_original_null || v < 0) { .. }
\end{lstlisting}

This example suggests that our approach can affect programs in ways not directly
related to I/O operations.
%
In C, programs can use a variable that stores a stream (e.g., for nullity
checking) even after closing it.
%
In contrast, Rust allows the use of only owned variables, and closed stream
variables are no longer considered owned.
%
As a result, any statements that access closed stream variables can be affected
by the transformation, potentially making the program uncompilable.
%
This implies that developers may need to slightly revise their code before
translation to satisfy Rust's strict ownership discipline.
%
In addition, although this specific case is not a bug, it supports the common
belief that translating to Rust can reveal previously unknown bugs, such as
using a stream for I/O operations after it has been closed.

The experimental results also show that our transformation likely preserves
program behavior.
%
All the 32 programs still pass the tests after the transformation.

To further investigate correctness, we conducted case studies.
%
To effectively evaluate correctness, we need to focus on the possible sources of
incorrect transformation that arise from the design of our approach.
%
An incorrect transformation can result from:
%
(1) incorrect origin and capability analysis or an incorrect decision on stream
types,
%
(2) incorrect transformation of stream construction, assignment, or function
call expressions, or
%
(3) incorrect error propagation.

Among these, the first two do not require manual checking.
%
First, an incorrect stream type results in a type error and is automatically
detected by the compiler.
%
Second, we transform stream construction, assignment, and user-defined function
call expressions only to match the types of the left and right values, and any
incorrect transformation also leads to a type error.
%
In addition, since the transformation of API function calls relies on the same
mappings across all programs, it is meaningless to perform case studies on
individual translations.

Therefore, our case studies focus on the correctness of error propagation.
%
Since each error-checking call is independently transformed, each case study
targets a single error-checking call.
%
Among the 62 programs, we found 140 transformed error-checking calls across 24
programs.
%
In each program, we studied all calls if it contained three or fewer, and
randomly sampled three otherwise, resulting in 59 calls studied in total.
%
The authors and another Rust expert independently reviewed each call, and we
confirmed that their assessments agreed.

Our studies reveal that the error propagation of 8 calls (1 from bc-1.07.1, 2
from gprolog-1.5.0, 1 from nettle-3.9, 2 from patch-2.7.6, 1 from time-1.9, and
1 from uucp-1.07) is potentially incorrect.
%
The reason is that these programs use error-checking calls to detect errors from
multiple API calls, rather than a single call, as illustrated below:

\begin{lstlisting}[language=C]
fputc('a', x); fputc('b', x); if (ferror(x)) { .. }
\end{lstlisting}

\noindent
%
Since our analysis considers only the API call found first while moving backward
as the source, earlier calls are not treated as sources, and the indicator value
is not propagated from them.
%
Therefore, the transformed code can behave incorrectly when an earlier call
fails but a later call succeeds.
%
We believe this issue is rarely problematic in practice, because if an earlier
call fails, a later call is likely to fail again.
%
Nevertheless, addressing this issue would be an interesting direction for future
work.

\subsection{RQ2: Efficiency}
\label{sec:evaluation:efficiency}

\begin{figure}[t]
  \centering
  \includegraphics[width=0.78\columnwidth]{imgs/time.pdf}
  \vspace{-1em}
  \caption{Forcrat execution time}
  \vspace{-1.5em}
  \label{fig:time}
\end{figure}

We evaluate the efficiency of the proposed approach by measuring the time taken
to analyze and transform each program.
%
The experimental results show that our approach is highly efficient.
%
\cref{fig:time} presents the origin and capability analysis time, error source
analysis time, transformation time, and the total time, all relative to the Rust
LOC.
%
Even for the largest program, gsl-2.7.1 with 422k LOC, the entire process takes
less than 14 seconds.
%
Note that the time required to translate each benchmark program using C2Rust
ranges from 0.6 to 295 seconds.
%
On average, running the proposed method on C2Rust's output takes only 8.9\% of
the time spent by C2Rust, demonstrating the efficiency of the proposed approach.

The results provide empirical evidence that the proposed approach scales well to
large programs in practice.
%
Although the origin and capability analysis has cubic time complexity in theory,
it exhibited nearly linear performance in our experiments.
%
This is because real-world programs contain a relatively small number of
stream-storing locations compared to their overall size.
%
The majority of analysis time is spent collecting constraints by visiting each
statement, rather than solving the constraints.
%
To confirm this, we also measured the constraint-solving time and found that it
completed in less than 5 milliseconds for each program.

Similarly, the error source analysis is quadratic in theory, but it also showed
nearly linear performance in our experiments.
%
This is because error-checking calls and their sources are typically located
close to each other, allowing the source to be found after analyzing only nearby
functions.
%
We measured the number of functions analyzed per error-checking call and found
that only 2.15 functions were analyzed on average.

\subsection{RQ3: Applicability}
\label{sec:evaluation:applicability}

We evaluate the applicability of the proposed approach by counting the number of
libc I/O API calls replaced in each program.
%
If a stream is classified as unsupported by the analysis, its type remains
\C{*mut\ FILE} during transformation, and the I/O API calls on that stream are
not replaced with Rust std API calls.
%
The practical value of our approach is high only if it can replace the majority
of API calls.
%
Therefore, we assess whether our approach is applicable to a sufficiently large
number of API calls.

The experimental results show that our approach achieves high applicability by
replacing 11,257 out of 13,594 libc I/O API calls (82\%) across all benchmark
programs.
%
The last column of \cref{tab:benchmark} shows the number of libc I/O API calls
that remain after the transformation.
%
Notably, in 28 out of the 62 programs, Forcrat successfully replaces all API
calls.
%
We believe these results are promising, especially considering the design
differences between C and Rust.
%
For example, \textsc{Crown}~\cite{zhang2023ownership} replaces 30\% of raw
pointers with references in a subset of these benchmark programs, highlighting
the inherent difficulty of translating C features into Rust.

Nevertheless, future research is needed to replace more API calls.
%
For example, Forcrat exhibits a low conversion ratio in a few programs:
%
below 50\% for four programs (gawk-5.2.2, libtool-2.4.7, rcs-5.10.1, and
sed-4.9).

\begin{figure}[t]
  \centering
  \includegraphics[width=0.9\columnwidth]{imgs/unsupported.pdf}
  \vspace{-1em}
  \caption{Reasons for failing to replace API calls}
  \vspace{-1.5em}
  \label{fig:reason}
\end{figure}

To guide future research directions, we investigate the reasons for failing to
replace each API call, which correspond to the reasons for classifying the
stream used as the argument of the call as unsupported.
%
We follow the 11 reasons described in \cref{sec:analysis:unsupported} in
this investigation.
%
Note that a single call can have multiple reasons, e.g., when the argument
stream is passed to \C{setbuf} and also used in a pointer comparison.
%
\cref{fig:reason} presents the results:
%
blue bars show the number of API calls affected by each reason, and red bars
show the number of calls affected \emph{solely} by that unique reason.

The results show that many calls are affected by \emph{multiple} reasons, so
resolving a single reason alone would not lead to significant improvement.
%
That said, Comparison and Casts appear as the unique reason for more than a
hundred calls, making them higher-priority targets for future work.
%
To address Comparison, it may be feasible to compare the underlying file
descriptors of streams, as Rust std supports obtaining file descriptors from
streams.
%
For Casts, since Wu and Demsky's work~\cite{wu2025genc2rust} replaces \C{void}
pointers with Rust generics, adopting their approach could handle such cases.
%
In addition, we may reduce the number of streams that are unsupported due to
Improper Capabilities by developing a flow-sensitive version of the origin and
capability analysis.
%
If a stream variable has different origins and capabilities at different program
locations, a flow-sensitive analysis can capture this information, allowing us
to split the variable into multiple variables with different types.

\subsection{RQ4: Impact on Performance}
\label{sec:evaluation:performance}

We evaluate the impact of replacing libc streams with Rust std streams on
program performance.
%
We first investigate how many \C{dyn}-type variables are introduced by the
transformation, as \C{dyn} incurs run-time overhead due to virtual dispatch.

It turns out that \C{dyn} variables are rarely introduced---only 0.9 per program
on average.
%
In 48 programs, no \C{dyn} variables are introduced at all.
%
The maximum occurs in cflow-1.7, where 14 \C{dyn} variables are introduced.
%
We believe this is because C programmers typically handle streams from different
origins by passing them to the same function, but rarely by storing them in the
same variable.

In addition, we compare the performance of each Rust program before and after
the transformation by measuring the execution time of the test suite.
%
To ensure reliable results, we used only the 25 programs whose test suite
execution time exceeds 0.1 seconds and ran each program fifty times.

The experimental results show that the performance difference is negligible.
%
The transformed programs were slower by only 1.23\% on average compared to the
original ones.
%
In 9 programs, execution was even faster after the transformation.
%
We also performed a one-sided Welch's t-test on the results from each program
with the null hypothesis that the transformed program is slower than the
original by 5\%.
%
With a significance level of 0.05, we could reject the null hypothesis for 15
out of 25 programs, implying that the transformation is unlikely to incur
significant performance overhead.
%
While we could not reject the null hypothesis for the remaining 10 programs,
this does not necessarily mean that the overhead exceeds 5\%.
%
Increasing the number of runs may yield stronger statistical significance,
allowing us to reject the null hypothesis.

\subsection{RQ5: Security}
\label{sec:evaluation:security}

We evaluate how much the proposed approach improves the security of programs.
%
In Rust, code security is often approximated by measuring the amount of
\emph{unsafe} code because the compiler cannot verify its safety and leaves the
burden of validation to the developers.
%
Following this convention, we measured the number of characters in the unsafe
code before and after the transformation.

The results show that the transformation slightly improves security by reducing
unsafe code by 4.4\%.
%
This reduction comes from replacing API calls, as all libc API calls are
considered unsafe by the Rust compiler, while most Rust std API calls are
considered safe.
%
Unfortunately, the security improvement is small because unsafe code arises from
various sources~\cite{emre2021translating}.
%
Achieving a significant improvement requires addressing diverse features and
APIs, and our work represents an important step toward this long-term goal.

% Note that prior studies~\cite{emre2021translating, emre2023aliasing,
% zhang2023ownership, wu2025genc2rust, hong2024dont, hong2024tag, hong2023concrat}
% did not measure unsafe code reduction, focusing only on the conversion rate of
% the targeted feature.

\subsection{RQ6: Code Changes}
\label{sec:evaluation:code}

We evaluate how much the transformation changes the code by counting the line
differences between the code before and after the transformation.
%
This helps estimate the effort required to manually replace the I/O API, and
also
provides insight into the verbosity of the transformed code.

The results show that the transformation relieves a significant burden but
increases code verbosity.
%
On average, 1,904 lines are inserted and 553 lines are deleted per program.
%
Manually modifying the code to this extent would require substantial effort.
%
At the same time, this indicates that each line is replaced with roughly four
lines.
%
This increase in code size is mainly due to
%
(1) some libc API functions being replaced with sequences of Rust std API calls,
owing to the lack of exact equivalents, and
%
(2) error propagation requiring the introduction and updates of error variables.
%
Reducing the verbosity would be an interesting direction for future research.

\subsection{Threats to Validity}
\label{sec:evaluation:threats}

Threats to external validity stem from the choice of benchmarks, which consist
of GNU packages and open-source projects on GitHub.
%
Since GNU packages often share common code patterns, including many GNU packages
in the benchmarks may introduce bias.
%
Although open-source projects enhance diversity, they still cannot fully
represent the entire C ecosystem.
%
Expanding the experiments to include a wider variety of C programs would
strengthen the evidence for the generalizability of our approach.

Threats to construct validity arise from using test suites for correctness.
%
Since tests do not cover all behaviors of the programs, passing all tests does
not guarantee semantics preservation.
%
However, we discovered bugs in our implementation during development, indicating
that the test suites are reasonably effective.
%
That said, testing the transformed programs more thoroughly using fuzzing or
automated unit test generation would provide stronger evidence of correctness.

Threats to construct validity arise also from using test suites for performance
comparison.
%
The test suites are not designed for performance measurement, and many exhibit
high variance in execution time.
%
To mitigate this, we ran each program multiple times and used statistical tests.
%
Furthermore, we discovered improper API mappings that led to performance
degradation during development, indicating that the test suites are suitable for
revealing performance issues.
%
However, developing dedicated benchmarks for performance evaluation would
facilitate a stronger analysis of performance impact.
